% SPDX-License-Identifier: CC-BY-4.0
% Copyright (C) 2026 the Langlandscape team
% Section "Agentic setup". Report body: anonymised (no person, machine, account or user names).
% Every number is taken from the repository file named next to it; nothing here is estimated.
\section{Agentic setup}\label{sec:agentic}

\begin{sloppypar}
All mathematics, Lean code, computations and talk material in this repository were produced by a hierarchy of
Claude Code agents, coordinated through git and working across four machines: \emph{the main orchestration machine},
two Lean build machines (\emph{build machine A}, a 16-thread workstation; \emph{build machine B}, a 16-thread workstation),
and \emph{the second team member's machine} (a notebook), which runs
its own Claude Code session under \emph{a second, separate Claude account} (hardware:
\pathtt{agents/rules/resources.md}, \pathtt{docs/agentic-process/architecture.md}~\S1). The two team members designed this setup,
steered it in real time and took the critical decisions recorded in \pathtt{agents/status/STATE.md} (``Decisions''; \S\ref{ss:human});
questions awaiting a team decision were applied provisionally and marked as such (\pathtt{docs/agentic-process/workflow.md}~\S1). The full
process write-up is \pathtt{docs/agentic-process/} (README, architecture, workflow, timeline, metrics, lessons,
model study, HACS). Unless stated otherwise, the counts below are those of
\pathtt{docs/agentic-process/metrics.md} and \texttt{timeline.md}, state 2026-11-14, 18:10\,PST, i.e.\ after stage~6
and before the final integration. All times PST (UTC$-$8). Source logs: all machines on Berlin time (CET, UTC+1), report sources in UTC; converted.
\end{sloppypar}

\subsection{Orchestration architecture}
The redacted session transcripts behind this section are in \pathtt{supplementary/agent-logs/} (see \S\ref{sec:ai}).

\begin{sloppypar}
\paragraph{Hierarchy.} A single long-running \emph{global master} session (Opus) on the main orchestration machine
ran the project from 2026-11-14, 01:20\,PST, within the setup designed by the team and under its steering (continuous at the critical decisions, not at every hour; \S\ref{ss:human}; \pathtt{docs/agentic-process/timeline.md}~\S2). Its first
principle is to keep its own context clean: it delegates \emph{every} task, including reading large files and running
builds, and acts itself only to plan, take operational decisions, brief, merge into \texttt{main} and refer the critical decisions to the humans
(\pathtt{agents/roles/global\_master.md}). Below it sit \emph{stage masters} (coordinators, Opus; one per planning
phase and per execution stage, plus two HACS lane masters and a final-integration master), which split their stage
into work packages (WPs) and delegate them to \emph{workers}; every coordinator follows the base role
\pathtt{agents/roles/coordinator.md}. The shape is
\begin{center}
\small
global master $\to$ stage masters $\to$ (track sub-coordinators $\to$) workers, auditors, experiment strands,
\end{center}
with Lean build jobs and the second machine attached to the global master only
(\pathtt{docs/agentic-process/architecture.md}~\S2). Each master and worker works in its own git worktree on its own
branch; subagents commit but never push; only the global master pushes and merges into \texttt{main}
(\pathtt{agents/rules/orchestration.md}). The concurrency cap is \textbf{8 concurrent Claude subagents per delegating
agent} (raised from 3 on 2026-11-14 at 07:32\,PST), subject to the memory and disk thresholds of
\pathtt{agents/rules/resources.md} (at least 12\,GB RAM kept free, about 2\,GB budgeted per subagent); Fable agents may
run in parallel with all others.

\paragraph{Models by task class.} The policy is \pathtt{agents/rules/model-selection.md}, summarised in
Table~\ref{tab:agentic-models}. It was not fixed in advance: it is the outcome of a model-versus-task study run inside
the stages (\pathtt{plan/implementation/EXPERIMENT\_PROTOCOL.md}, \pathtt{docs/agentic-process/model-study.md}).
For a selected WP, two (in stage~1 up to three) models implemented the same brief from the same base commit at the
same minute; build machines and auditors saw only neutral job ids, never model names; each strand got its own blind
auditor, a rubric score and a weighted finding count $W = 3\cdot\mathrm{crit} + 2\cdot\mathrm{maj} + \mathrm{min}$;
the decision metric was the active time (implementation $+$ audit $+$ fix minutes): the fastest passing strand wins,
unless another passing strand of equal quality is within $\tau = \max(15\,\%\mbox{ of the best}, 5\,\mathrm{min})$,
in which case the cheaper one wins. Twelve such comparisons were run in stages 1--5
(25 local strand branches; \texttt{metrics.md}~\S1). Examples of the evidence behind the table: on routine Lean,
Sonnet passed in 4.0 active minutes against 13.0 for Muse Spark Open, which failed with two critical findings
(stage~1); on routine LaTeX an Opus single shot (7.8 min) failed on the minor count while Sonnet plus one Opus fix
(8.05 min) passed (stage~2); on hard Lean off the critical path, Sonnet plus one Opus fix took 18.1 active minutes
against 26.3 for Opus at less than half the cost index (stage~4); on hard LaTeX, Opus with a planned fix round
passed while a Fable single shot failed on the minor count (stage~3). Every verdict is a prior from one or two
comparisons per cell (\pathtt{model-selection.md}~\S2; confounds such as remote build latency and Opus auditing
Opus are listed in \pathtt{model-study.md}~\S4).
\end{sloppypar}

\begin{table}[ht]
\centering\small\sloppy
\begin{tabular}{@{}p{0.38\linewidth}p{0.22\linewidth}p{0.32\linewidth}@{}}
\toprule
Task class & Model & Fix-round budget \\
\midrule
Lean mechanical (copying, vendoring) & Haiku 5.5 & none; red check $\to$ Sonnet \\
Lean routine, Lean ``mechanical with meaning'' & Sonnet 5.5 & one Opus fix \\
Lean hard on the critical path & Opus 5.5 (right first time) & one planned Opus fix \\
Lean hard off the critical path & Sonnet first & one Opus fix \\
LaTeX and code, routine & Sonnet 5.5 & one Opus fix (6--8 minor findings expected at round~0) \\
LaTeX hard (full proofs) & Opus 5.5 & one planned Opus fix, always \\
conceptual, delicate; Opus failures & Opus or Fable 5.2 & --- \\
Lean audits, chapter audits & Fable 5.2, in parallel & --- \\
coordination (stage masters) & Opus 5.5 & --- \\
logging, polling, axiom lists, file moves & Haiku 5.5 & --- \\
opencode (Mistral Large 4, Muse Spark Open) & Q\&A only & not used as implementers \\
\bottomrule
\end{tabular}
\caption{Model policy by task class (\pathtt{agents/rules/model-selection.md}~\S1). A fix that does not reach green
escalates one rung (Haiku $\to$ Sonnet $\to$ Opus $\to$ stage master, who splits, re-scopes or calls Fable); Fable is
never on an escalation ladder.}\label{tab:agentic-models}
\end{table}

\begin{sloppypar}
\paragraph{Audits.} Every stage closed with Fable audits run in parallel on the early-gate commit: a Lean audit
(\pathtt{agents/roles/lean\_audit.md}) and a chapter audit (\pathtt{agents/roles/mathematics\_chapter\_audit.md});
from stage~4 on, both were run in parts while WPs were still open, and their required remarks became fix WPs inside the
same stage (\pathtt{docs/agentic-process/workflow.md}~\S\S3, 5). Opus ran the WP-level audits. The Fable audits found
what the Opus audits had missed: the only discharge of \pathtt{Hyp.RamNrdSurj} in stage~2 was vacuous (it covered
$D = 1$ only); in stage~6 the registered hypothesis \pathtt{Hyp.FreeRankTwo} turned out to be equivalent to the
Galois half of the conclusion of \pathtt{Defect.defect\_formula}, which is now disclosed as such everywhere the theorem
is stated; three of four ``non-vacuity reasons'' given by writers were false; and the stage-6 chapter failed for
$M > 1$ before it was restricted (\pathtt{docs/agentic-process/lessons.md}, items 11, 32--34;
\pathtt{model-study.md}~\S3a). The per-stage findings are tabulated in \texttt{metrics.md}~\S6. The consolidated
final report that this report builds on was audited by Fable against the chapters, the result files and the gate logs
(\pathtt{mathematics/audits/2026-11-15\_final-report-audit.md}: FAIL with three blockers, all one-line text
corrections, PASS WITH REMARKS once they are applied; all blockers, majors and minors were applied in one fix round, merge \texttt{b2f6aa20}, and a cross-account review followed with a second fix round).

\paragraph{Cross-account reviews.} Reviews by the same account share its blind spots, so finished chapters were also
read \emph{blind} by agents of the second, separate Claude account on the second team member's machine: the chapters
of stages 1--5 (\pathtt{mathematics/audits/2026-11-1*\_s0*-xacct.md}), second readers for conceptual claims of stage~6,
and the final report and talk (\pathtt{mathematics/audits/2026-11-15\_final-xacct.md}: verdict FIX LIST, 0 critical,
7 major, 18 minor). This independent reading changed a statement of the mathematics twice: a gap in the
$q \equiv -1$ claim for $M > 1$ was found there, the $M = 1$ case confirmed, and a counterexample restricted the
statement (\texttt{lessons.md}, item 34). The second machine also wrote WPs (stage-6 commutative-algebra Lean files,
chapter drafts, literature checks); its Lean was accepted only after a build on build machine A or B.

\paragraph{The message protocol between machines.} Native session messages reach only sessions of the same account.
Between the two accounts all traffic is files in git (\pathtt{agents/rules/messages.md}): each machine has an inbox
directory under \pathtt{agents/messages/}; a message is a Markdown file named
\texttt{YYYYMMDD-HHMMSSZ\_<from>\_<slug>.md} with a front matter (id, from, to, type
$\in$ \{request, report, question, answer, ack, info\}, re, priority), committed and pushed on \texttt{main}; the
recipient answers or acknowledges and moves the file to \texttt{archive/}, so an inbox root holds exactly the
unprocessed messages. Inboxes were polled every 10 minutes. By 18:10\,PST on 2026-11-14 there were 134 message files
and 7 work batches had been sent (\texttt{metrics.md}~\S\S1, 5); batch~8 (review of the final report and talk) followed.
Inside the main machine, stage masters talk to the global master through committed \emph{outbox} files
\pathtt{plan/implementation/outbox/<UTC>\_<id>.md} (job specifications, status, milestone closures) and \emph{inbox}
files with verbatim build results; a session message is only a notification, never the payload. This rule (decision
F11, 2026-11-14 08:49\,PST) followed an incident in which in-session messages were delayed by hours and one job batch
was lost; every channel that writes to git was reliable (\texttt{lessons.md}, item 5). Build results come back on
result branches named after the build machine and the job id, which are read but never merged; deliveries of the second
machine come on branches of their own and are merged by the stage master after acceptance
(\pathtt{architecture.md}~\S4).

\paragraph{opencode.} Two non-Claude models were available through the opencode CLI
(\pathtt{agents/rules/opencode.md}): Mistral Large 4 and Muse Spark Open. In the planning phases Muse Spark Open served as
a cheap first-pass auditor behind Claude revisers who judged every finding (7 of 70 findings confirmed, 22 partially,
41 rejected in the staged plan; 15 of 89 confirmed in the implementation plans, where 12 findings labelled critical had
none left after verification). Both models competed as implementers in the stage-1 model study and won no comparison
(\pathtt{model-study.md}~\S2). From then on they were used for
questions and answers only, never as implementers (\pathtt{plan/implementation/OVERVIEW.md}~\S5.1, decision F13).
\emph{One disclosed exception:} the discharge of \pathtt{Hyp.RamNrdSurj} was drafted by an external model
(\pathtt{mistral-large-4-0} via opencode), run by a team member on the other team member's machine, outside our agent
process. It was then integrated into the library by the agents as \pathtt{Langlandscape.Hyp.ramNrdSurj\_holds}, audited
independently by Fable (PASS WITH REMARKS), and gated from scratch (build machine A, \pathtt{ramnrd-final-2}, 840/840
standard axiom lines; \pathtt{docs/agentic-process/workflow.md}~\S6). The gate is listed in
Appendix~\ref{sec:commits}.
\end{sloppypar}

\subsection{The wall-time minimizer}

\begin{sloppypar}
The design goal of every rule above is to minimise \emph{wall} time, not agent time or cost
(\pathtt{agents/rules/orchestration.md}). The mechanisms, each recorded in the stage run files
\pathtt{plan/implementation/runs/s0*.md}:
\begin{itemize}
\item \emph{Pipelining.} Tracks are decoupled: mathematics may run ahead of Lean, and no milestone of one track waits
for another; the only coupling is that a statement must exist in the mathematics before it is formalized
(\pathtt{orchestration.md}, ``Decoupled tracks''). Lean WPs stack on unverified drafts of their upstream WPs instead of
waiting for them (rule F14). Stages overlapped: the stage-6 commutative-algebra layer was written on the second machine
during stages 2--3 and merged in stage~3's first integration round (\texttt{runs/s03.md}); each finished chapter went
to its Fable audit and cross-account review as soon as it existed; the two HACS lanes ran on the main machine's GPU
next to stages 5 and 6 (\pathtt{docs/agentic-process/hacs.md}).
\item \emph{Early gates.} An early full gate (gate-1) runs as soon as the core Lean WPs are green, and from stage~5 a
\emph{pregate} on the merged tree without the open WPs at $t \approx 35$\,min (709/709 in stage~5, 786/786 in stage~6),
so that gate-1 became a short round (6 and 7 minutes; \texttt{workflow.md}~\S5). In stage~5, gate-1 passed at
$t = 43$\,min against a planned $\approx 80$\,min (milestone closure s05).
\item \emph{Model policy and fix-round budgets} (Table~\ref{tab:agentic-models}): one planned fix round per WP is
cheaper than a stronger first shot; after two failed Sonnet rounds a WP escalates to Opus, and a parallel rescue writer
starts at round~3 (\texttt{lessons.md}, item 20).
\item \emph{Parallel Lean job distribution.} Batched integration rounds amortise the remote round trip; from stage~3 two
build machines ran single-WP lettered jobs in parallel. Stage~3 ran 22 rounds in 52 minutes with 2--3 minutes per round
(\texttt{runs/s03.md}, lesson 1); stage~4 ran 53 job ids in 112 minutes, typically 2--4 minutes per result
(\texttt{runs/s04.md}, lessons 8--9).
\item \emph{Warm seeds.} Each build machine keeps a warm-seed worktree, advanced after each green job to the latest
verified integration point, and seeds new jobs from it by copy-on-write
(\pathtt{agents/roles/lean\_build\_node.md}, ``Warm seed''). Seeded versus from-scratch builds of the same tree:
stage-5 gate-1 223\,s vs 279\,s; stage-6 gate-1 20\,s vs 304\,s (milestone closures s05, s06). Final gates are never
seeded.
\end{itemize}

Table~\ref{tab:agentic-walls} gives the result. The stage wall is measured from the stage's $t_0$ to its final-gate
PASS. Stages 1--6 together took $\approx 10.4$\,h including a 1\,h re-plan after stage~1, against a plan of
$\approx 21.7$\,h; the planned Lean walls were repeatedly about twice too long, and the critical path was a chain of
cheap single-error fix rounds rather than proof difficulty (\texttt{timeline.md}~\S1, \texttt{runs/s04.md},
\texttt{runs/s06.md}).
\end{sloppypar}

\begin{table}[ht]
\centering\small\sloppy
\begin{tabular}{@{}lllll@{}}
\toprule
Stage & Wall (h) & Plan (h) & Actual/plan & Source \\
\midrule
1 quaternion foundations & 2.2 & 8.8 & 0.25 & \texttt{runs/s01.md} \\
2 tree quotients & 1.3 & 2.5 $\pm$ 0.7 & 0.52 & \texttt{runs/s02.md} \\
3 Ribet--Ihara & 1.1 & 2.85 (2.2--3.6) & 0.39 & \texttt{runs/s03.md} \\
4 $\Lambda$-adic Ribet & 2.2 & 3.0 (2.6--3.8) & 0.73 & \texttt{runs/s04.md} \\
5 comparison & 1.1 & 2.1 (1.5--3.2) & 0.52 & \texttt{runs/s05.md} \\
6 integral defect & $\approx$ 1.2 & 2.2--2.6 & $\approx$ 0.5 & \texttt{runs/s06.md} \\
\bottomrule
\end{tabular}
\caption{Stage wall times against plan. The planning phases before stage~1 (staged plan, Fable seam audit,
implementation plans, Fable plan audit) took about 6\,h. Source: \pathtt{docs/agentic-process/timeline.md}~\S1.}\label{tab:agentic-walls}
\end{table}

\subsection{Lean across machines}

\begin{sloppypar}
\paragraph{Frozen toolchain.} Lean \texttt{v4.35.0} and Mathlib at revision
\texttt{56a005855ea17e0b0958f73ec801ae140b0a1679} are frozen as a local copy on every machine
(\pathtt{agents/rules/frozen.md}). \texttt{lake update}, \texttt{lake clean}, \texttt{lake exe cache get}, any edit of a
toolchain or manifest file and any write into the frozen tree are forbidden for all agents. Builds run only through
wrapper scripts, which take locks, cap memory, check a pin file of SHA-256 sums of the manifest files and abort with exit
code 99 on any mismatch; \pathtt{check\_frozen.sh}, part of every gate, verifies that the frozen trees are untouched. It
caught one bad import (a Mathlib module that does not exist at the frozen revision) before it could do harm
(\texttt{workflow.md}~\S4, \texttt{lessons.md}, item 13).

\paragraph{Job flow.} From about 07:50\,PST on 2026-11-14, all Lean builds ran on the two build machines, because the main
machine's CPU and GPU were reserved for the HACS computations (\pathtt{agents/rules/resources.md}, ``TEMPORARY''). A job
goes through four steps: (1) the stage master commits an outbox file with the job specification (which tree, which
targets, which checks, the expected number of axiom lines); (2) the global master pushes the needed branch (for
experiment strands a neutral temporary ref) and relays the specification to a build machine with a free slot (machine A
has three slots, machine B two); (3) the build machine builds through the wrappers and pushes its logs to a result
branch for that job; (4) the global master writes an inbox file with the verbatim errors and goal states (never fix
hints) and the stage master integrates: a WP is merged only after its check passed \emph{on the merged tree}
(\pathtt{architecture.md}~\S\S1a, 3; \texttt{workflow.md}~\S3). The second machine's Lean drafts entered only through
acceptance jobs on build machine B.

\paragraph{Gates.} A gate has the steps (a) build of the library and its tests, (b) build of the vendored stacks,
(c) \pathtt{Audit/PrintAxioms.lean} with a count of lines against the expected number, every line a subset of
\{\texttt{propext}, \pathtt{Classical.choice}, \texttt{Quot.sound}\}, (d) a forbidden-token grep, (e) a
\texttt{sorry} count, (f) \pathtt{check\_frozen.sh}, plus per-file checks of the stage's files. Every stage's final gate
was built \emph{from scratch} (unseeded) on build machine A; from stage~2 on it was cross-checked on build machine B
(build machine B came online at 11:48\,PST on 2026-11-14, after stage~1 had closed; \texttt{timeline.md}~\S2). All six
final gates passed with standard axioms only and no \texttt{sorry} in our code, with 161, 195, 519, 592, 747 and 836
axiom lines; the gate after the \pathtt{Hyp.RamNrdSurj} discharge (job \pathtt{ramnrd-final-2}, build
machine A, from scratch) built 4198 jobs in 321\,s and printed 840/840 standard lines, with build machine B's quick check
\texttt{ramnrd-3} also passing; the gate after the Lean submission bundle (job \texttt{subl-gate-1}, build machine A, from
scratch, Lean tree 28d2b091), built 4199 jobs in 323\,s and printed 950/950 standard lines, with build machine B's fresh-checkout
cross-check \pathtt{subl-gate-1x} also passing (\texttt{metrics.md}~\S3; \texttt{timeline.md}~\S3); and the final gate, after the merge of
the \pathtt{Hyp.BrandtStabFree} discharge (job \texttt{bsf-gate-1}, build machine A, from scratch in an unseeded worktree, Lean tree
c70ef3a9), built 4205 jobs (and the vendored files, 1955 jobs) and printed 955/955 standard lines. The only \texttt{sorry} in the built
tree lies in the frozen import (\pathtt{AlgJL/VolIII/P2/Instances.lean}), outside every audited declaration.

\paragraph{Job counts.} Table~\ref{tab:agentic-jobs} counts the job traffic per stage, as outbox and inbox files and as
result branches per build machine (\texttt{timeline.md}~\S3, \texttt{metrics.md}~\S7). In total 107 result branches of
build machine A and 68 of build machine B existed at 18:10\,PST on 2026-11-14, among them acceptance jobs for the second
machine's Lean.
\end{sloppypar}

\begin{table}[ht]
\centering\small\sloppy
\begin{tabular}{@{}lrrrrp{0.36\linewidth}@{}}
\toprule
Stage & Outbox & Inbox & Res.\ A & Res.\ B & Final gate: axiom lines; from-scratch wall \\
\midrule
1 & 19 & 0 & 16 & 0 & 161; 3842 jobs \\
2 & 39 & 16 & 19 & 1 & 195; cross-check B PASS \\
3 & 60 & 30 & 20 & 11 & 519; A 180\,s, B 3:16 \\
4 & 85 & 55 & 30 & 25 & 592; A 251\,s, B 3:50 \\
5 & 42 & 24 & 11 & 13 & 747; A 273\,s, B 4:22 \\
6 & 47 & 21 & 10 & 13 & 836; A 319\,s, B 5:21; 4194 jobs \\
\bottomrule
\end{tabular}
\caption{Lean job traffic per stage. Outbox/inbox: files whose name contains the stage (a proxy that also counts status
and closure files; the inbox was introduced in stage~2). Results A/B: result branches of build machines A and B.
Sources: \pathtt{docs/agentic-process/timeline.md}~\S3, \texttt{metrics.md}~\S\S3, 7, milestone closures
\pathtt{plan/implementation/outbox/*milestone-closure-s0*.md}.}\label{tab:agentic-jobs}
\end{table}

\begin{sloppypar}
\paragraph{Scale.} At 18:10\,PST on 2026-11-14, \texttt{main} had 2058 commits (405 merges); the Lean library had 195
non-vendored files with about 37\,980 lines (\texttt{metrics.md}~\S\S1, 3); 29 hypotheses were registered in
\pathtt{lean/Langlandscape/HYPOTHESES.md}, of which 2 were then discharged and 1 conditionally discharged (now 3 and 1). Token counts and
costs per model are not recorded in the repository.
\end{sloppypar}
